\chapter{7}

\section{Mailand (IT), Montag, 20.
August}

Mit gebeugten Fingern strich er ihr sanft über die Wange. Sie war schön
anzusehen. Weit entspannt lag sie in seinem Arm und schlief.
Vertrauensvoll wie ein Kind.

Seine kostbare Liebe. Fast schon eine junge Frau.

«Amore, es ist Zeit zu gehen,» sagte er mit wacklig zärtlicher Stimme.
Er nahm einen Schluck Wasser aus einer Trinkflasche. Es fiel ihm schwer,
seine Bewegungen zu koordinieren. Er musste sich zusammenreissen.

Er legte die Lederjacke etwas dichter an ihren Hals. Es war kühl
geworden. Die Luft feucht.

Das Mädchen öffnete seine Augen zu schmalen Schlitzen. Schloss sie
wieder, blinzelte und sprach kaum verständlich vor sich hin. Etwas von
schön und Übelkeit, etwas von Angst und himmlischen Gefühlen, von Gehen
und von ihrer Mutter…er verstand nicht.

Auch dieses Mal hatte es ihn Überwindung gekostet. Sie hatte darauf
bestanden, mitzumachen. Er schliesslich nachgegeben, die Dosis halbiert.
Wenigstens, hatte er sich gesagt.

Wieder lagen sie hier. An diesem magischen Ort. Geschützt von einem
hohen, starken Felsen. Der seitlich mit sommertrockenen Büschen
eingerahmte Wiesenplatz gehörte ihnen alleine. Mit offenem Blick nach
oben. Entlang der Baumstämme, bis hinauf zu ihren weiten Kronen.

Mit ihrem jungen Alter haderte er. Wenn sie sich zusammen auflösten und
entschwanden, vergass er es. Und schliesslich hatte sie ihn, charmant
zwar, aber unerbittlich, gedrängt.

«Ich will es ausprobieren, vertraue dir, als Profi. Ich bin kein Kind
mehr, ich weiss, was ich will.»

Dabei hatte sie ihn am Arm gepackt.

Dass sie noch im Schutzalter sei, liess sie nicht gelten. Sie setzte
jedes Mal einen trotzigen Blick auf. Sie gefiel ihm in ihrer
Lebendigkeit.

Dann war er eingeknickt und redete sich ein, er würde auf sie aufpassen.

Und jetzt musste er sie unverzüglich zu ihrer Freundin nachhause
bringen. Die wusste um ihre Beziehung. Als Einzige. Vorgelogene
Übernachtungen bei ihr waren Teil des Spiels geworden.

«Die Hand hat einen Arm,» flüsterte sie aufgewühlt mit unterdessen weit
aufgesperrten Augen.

Sie hatte ihre Visionen auch in dieser Nacht ohne Punkt und Komma
kommentiert. Mal lachend, weinend, schreiend, schluchzend. Mal staunend.
Dabei war sie weit entrückt und zwischendurch zitterte sie.

Ihre phantastischen Bilder beeindruckten ihn. Nachträglich würde er sich
schuldig fühlen.

«Einen nackten, einen blossen Arm,» schob sie langsam nach.

Sie hob ihre Stimme. Zog die Worte in die Breite. Es hörte sich
zerrissen an. Sie streckte ihren rechten Arm unsicher hoch und zeigte
mit dem Finger nach oben.

Er schaute sie ununterbrochen an.

\emph{«Carissima, ti voglio bene»},\footnote{«Meine Liebste, ich liebe dich.»} flüsterte er und näherte sich ihrem
Gesicht.

Sie stiess ihn hart zurück.

«Hör auf. Siehst du’s nicht? Ein Arm mit Hand, ganz echt, da auf dem
Baum, da, vielleicht ist noch mehr dran?»

Sie klang schockiert und blickte ihn ruckartig an. Ihre vergrösserten
Pupillen liessen ihre Augen wie tiefdunkle Kugeln aussehen.

Er, mit einem Mal nüchtern, starrte nun auch hinauf in die Baumkrone der
Akazie. Tatsächlich.

Zwischen den Blättern, am Rande, hing etwas hinunter. Etwas Helles.
Bewegungslos.

Ungelenk stand sie langsam auf. Er auch.

Es kostete ihn Kraft, einigermassen gerade zu stehen. Sie hielt sich an
seiner Hüfte fest. Wie Betrunkene stützten sie sich gegenseitig
aneinander ab. Und für einen langen Moment starrten sie ergriffen
fixiert und stumm auf das unwirkliche Bild.

Auch er war nun sicher. Dort oben hing ein Mensch fest. Oder auch nur
ein Teil davon. Er wollte sich darum kümmern. Anschliessend.

Erst musste er sie wegbringen. Nur keinen zusätzlichen Ärger
provozieren, hämmerte es in seinem Kopf.

Langsam tastete er nach ihrer Hand und zog sie zu ihren Taschen am
Boden. Sie liess es geschehen. Knickte immer wieder mit den Beinen ein
und drohte umzufallen.

Er redete leise, so kontrolliert wie möglich, auf sie ein. Sie staunte
ihn stumm an.

«Alles halb so wild, Amore mio. Du hast recht, das sieht aus wie ein
Menschenarm, aber, das ist keiner,» dazu liess er kurz ihre Hand los,
nahm schnell ihre Taschen hoch und bemühte sich, zu lachen, «ich glaub’s
ja nicht, eine \emph{Gummisusi,} schräg, wie die da hängt.»

Sie schaute ungläubig zu ihm hoch und wollte etwas sagen. Gut, dachte
er, dass sie nicht ganz bei sich ist. Er blieb stehen, umarmte sie mit
resolutem Griff und küsste sie auf den Mund.

Er wollte schnell vorankommen und keine Zeit verlieren. Er hob sie auf
seine Arme. Es war anstrengend. Seine Kraft war reduziert.

«Ich bringe dich zum Schlafen, meine Süsse,» säuselte er ihr ins Ohr.

Sie entspannte sich, und er trug sie mit hastigen Schritten über die
Wiese zu seinem Wagen.

Es war fast drei.
